\section{The \bench Benchmark}
\label{sec:benchmark}

\bench turns public attack recordings into investigation tasks in the form defined in \S\ref{sec:problem}.
Fig.~\ref{fig:pipeline} shows the construction; Table~\ref{tab:tasks} lists the tasks.

% Benchmark construction pipeline (TikZ, vector). Included by sections/benchmark.tex.
\begin{figure}[t]
\centering
\begin{tikzpicture}[
  font=\scriptsize,
  >={Latex[length=1.6mm,width=1.3mm]},
  bx/.style={draw=cA!70!black, fill=cA!9, rounded corners=1.5pt, align=center,
             text width=#1, minimum height=10mm, inner sep=1.5pt, line width=0.5pt},
  bx/.default=24mm,
  src/.style={draw=black!65, fill=black!6, rounded corners=1.5pt, align=center,
              text width=22mm, minimum height=7mm, inner sep=1.5pt, line width=0.5pt},
  res/.style={draw=cB!75!black, fill=cB!9, rounded corners=1.5pt, align=center,
              text width=25mm, minimum height=7.5mm, inner sep=1.5pt, line width=0.5pt},
  arr/.style={->, line width=0.7pt, black!70}]
% row 1: public recordings -> parse -> graph
\node[src] (o) at (1.2,0.45) {OTRF datasets\\[-1pt]Windows Sysmon};
\node[src] (s) at (1.2,-0.45) {Splunk attack\_data\\[-1pt]Sysmon for Linux};
\node[bx=25mm] (p) at (4.05,0) {\textbf{Parse, filter}\\JSON, XML (no DTD);\\drop self-telemetry};
\node[bx] (g) at (6.95,0) {\textbf{Graph builder}\\process, file, socket,\\account, host};
% row 2: selectors -> tasks / answers -> cue variants
\node[bx] (sel) at (6.95,-1.9) {\textbf{Scenario selectors}\\alert, core set,\\acceptable, anchors};
\node[res] (t) at (4.05,-1.45) {\texttt{tasks/} (agent)\\[-1pt]alert + graph};
\node[res] (a) at (4.05,-2.35) {\texttt{answers/} (scorer)\\[-1pt]ground truth};
\node[bx=22mm, draw=black!65, fill=white, densely dashed] (v) at (1.2,-1.45)
  {\textbf{Cue variants}\\\emph{blind}, \emph{mimicry}:\\strings renamed};
\draw[arr] (o.east) -- (p.west |- o.east);
\draw[arr] (s.east) -- (p.west |- s.east);
\draw[arr] (p) -- (g);
\draw[arr] (g) -- (sel);
\draw[arr] (sel.west |- t.east) -- (t.east);
\draw[arr] (sel.west |- a.east) -- (a.east);
\draw[arr] (t) -- (v.east |- t.west);
\end{tikzpicture}
\caption{Construction of \bench from public recordings. Agents see only \texttt{tasks/}; the scorer uses
\texttt{answers/}. The cue-controlled variants rename strings in the tasks and keep graphs and answers.}
\label{fig:pipeline}
\end{figure}


\subsection{Sources and selection}
We use two practitioner corpora recorded with real endpoint sensors. OTRF Security-Datasets~\cite{otrf_securitydatasets}
(MIT licence) provides Windows Sysmon~\cite{microsoft_sysmon} and Security event logs captured while ATT\&CK techniques were executed
in a lab domain; Splunk attack\_data~\cite{splunk_attackdata} (Apache-2.0) provides, among other formats, Sysmon
for Linux~\cite{microsoft_sysmonforlinux} recordings from Ubuntu hosts in Splunk's Attack Range. We selected recordings
that (i) contain a complete process lineage for the activity, (ii) are small enough to package (each source
file is at most 34\,MB), and (iii) together cover distinct tactics: credential access, persistence, privilege
escalation, defence evasion, execution, impact and service stopping. We excluded recordings whose key step is
not observable in the telemetry (for instance, an authorised-keys write that the sensor did not capture) and
very large recordings. The result is \BenchWinMal{} malicious Windows and \BenchLinMal{} malicious Linux tasks
covering \BenchTechniques{} distinct ATT\&CK techniques, and \BenchWinBen{}+\BenchLinBen{} benign tasks.

\subsection{Graph construction}
A graph builder maps events to entities and relations using the runtime's type vocabulary
(Table~\ref{tab:types}). Sysmon process-creation events create a process entity keyed by its process GUID, a
parent entity from the parent fields, a process-creation relation and an executed-file relation; network,
image-load, process-access, file-create, file-delete, registry, pipe, DNS and WMI events add the corresponding
relations. Windows Security events add scheduled tasks and logons, and service-installation events add
services. Linux events use the Linux type names (\tool{linux\_process}, \tool{linux\_file}, \tool{linux\_user},
\tool{linux\_process\_create}). Users are linked to processes, processes to hosts, and repeated identical
relations (for example a beacon reconnecting from new ephemeral ports) are collapsed with a count and a
last-seen timestamp. Field values are copied from the events, capped at 600 characters. Event files are
treated as untrusted input: XML lines declaring a document type or entities are rejected.

\paragraph{Collector self-telemetry} The Linux recordings include the log forwarder's own checkpoint-file
writes under its installation directory, thousands per minute. These are artefacts of the collection pipeline,
not host activity, so file-create and file-delete events under that directory are dropped. The forwarder's
processes and network connections are kept as benign background. For the five recordings where this churn
dominated, it reduced graphs from roughly 11{,}000--33{,}500 entities to 50--150, leaving the real host
activity.

% Entity and relation types across the benchmark (generated).
\begin{table}[t]\centering\scriptsize
\caption{Entity and relation types across \bench (counts over all tasks).}
\label{tab:types}
\begin{tabular}{@{}lr@{\hspace{1.2em}}lr@{}}\toprule
Entity type & Count & Relation type & Count\\\midrule
\texttt{linux\_process} & 3338 & \texttt{runs\_process} & 3881\\
\texttt{windows\_file} & 1034 & \texttt{executed\_file} & 3176\\
\texttt{linux\_file} & 1006 & \texttt{linux\_process\_create} & 3147\\
\texttt{windows\_process} & 543 & \texttt{loaded\_image} & 1837\\
\texttt{registry\_key} & 361 & \texttt{accessed\_process} & 726\\
\texttt{network\_connection} & 297 & \texttt{deleted\_file} & 704\\
\texttt{linux\_user} & 42 & \texttt{created\_registry\_key} & 282\\
\texttt{host} & 23 & \texttt{connected\_to} & 205\\
\texttt{windows\_user} & 23 & \texttt{wrote\_file} & 142\\
\texttt{named\_pipe} & 7 & \texttt{set\_registry\_value} & 137\\
 & & \texttt{accepted\_connection\_from} & 122\\
 & & \texttt{logged\_on\_to} & 70\\
 & & \texttt{win\_process\_create} & 29\\
 & & \texttt{resolves\_to} & 7\\
 & & \texttt{connected\_to\_pipe} & 7\\
 & & \texttt{deleted\_registry\_key} & 4\\
\bottomrule\end{tabular}\end{table}


\subsection{Scenario design}
\paragraph{Malicious tasks} Each technique recording contains an operator's shell session (a Windows
\texttt{cmd.exe} or PowerShell, a Linux login shell) from which the technique was executed. The alert is what
a detection would raise on the key step, for example ``Process opened LSASS with full access'' for a
\texttt{rundll32} invoking the MiniDump export of \texttt{comsvcs.dll}, or ``Kernel module loaded with insmod''.
Alert texts describe the observable event without stating a verdict.

\paragraph{Benign tasks} Benign tasks are real background events from the same recordings that a naive
rule would flag. They include \texttt{svchost.exe} opening LSASS with access mask \texttt{0x1000}
(query-limited information, insufficient to read memory) next to the malicious dump that opened it with
\texttt{0x1FFFFF}; the Windows Time service writing \texttt{\textbackslash RunTime\textbackslash} values that
match a careless \texttt{\textbackslash Run} rule; and message-of-the-day scripts running discovery commands
as root at SSH login. Two pairs are \emph{twins}: in the Linux service recording, an operator stopping and
disabling a web service and the package manager starting it during installation raise the same alert
(``Service state changed by systemctl as root'') with opposite verdicts.

\subsection{Ground truth}
Ground truth is written as selectors over the graph (image name, command-line substring, path, user,
destination address, relation type) rather than as identifiers, and is rebuilt from the raw recordings on
every build. For a malicious task, the \emph{session root} is the operator shell above the trigger (its parent,
or the nearest ancestor with a given image). The core set $C$ contains the trigger and the session root, the
key commands, the artefacts that carry persistence or impact (the dump file, the Run-key value, the module
copied into \texttt{/lib/modules}, the script that wiped directories), the host, and the account to act on:
the session owner, plus the trigger's account unless it is the system account the session elevated to
(\texttt{SYSTEM}, \texttt{root}). The acceptable set $A^\star$ adds every process in the session and every
file, registry key, connection and process the session touched. For a benign task, the anchors are the trigger
and the entities that prove legitimacy, for example the parent that launched it (the package manager, the
login daemon) or the relation carrying the decisive property (the access mask), each with its real property
fields.

\subsection{Statistics and format}
Graphs have \BenchEntMin{}--\BenchEntMax{} entities (median \BenchEntMedian{}) and
\BenchRelMin{}--\BenchRelMax{} relations (median \BenchRelMedian{}), with \BenchEntTypes{} entity types and
\BenchRelTypes{} relation types overall. The benchmark is a single \BenchZipMB\,MB zip: a manifest, one task
file per task (alert and graph, no answers), one answer file per task, and the upstream licences. A local
runtime serves the tasks through the same client interface as the competition runtime and scores
submissions, so an agent runs on \bench unchanged.

\begin{table*}[t]\centering\scriptsize
\caption{The 21 tasks of \bench. Graph size is entities/relations; \emph{Core} is the number of core entities (malicious) or legitimacy anchors (benign) in the ground truth.}
\label{tab:tasks}
\begin{tabular}{@{}lllllrr@{}}\toprule
Task & OS & Verdict & ATT\&CK & Alert title & Graph & Core\\\midrule
\texttt{lsass-dump-comsvcs} & Win & malicious & T1003.001 & Process opened LSASS with full access & 88/115 & 5\\
\texttt{logon-script-persistence} & Win & malicious & T1037.001 & User environment registry value written by reg.exe & 47/75 & 6\\
\texttt{service-binpath-hijack} & Win & malicious & T1543.003 & Service configuration changed with sc.exe & 157/219 & 6\\
\texttt{hta-startup-folder} & Win & malicious & T1218.005 & mshta.exe executed an HTA file & 338/717 & 6\\
\texttt{firewall-rule-added} & Win & malicious & T1562.004 & Inbound firewall rule added from the command line & 174/365 & 4\\
\texttt{python-http-server} & Win & malicious & T1059 & Interpreter from a user profile started a network service & 300/633 & 6\\
\texttt{bits-download} & Win & malicious & T1197 & BITS transfer job started from the command line & 186/247 & 4\\
\texttt{run-key-and-beacon} & Win & malicious & T1547.001 & Autostart Run key value set by PowerShell & 74/87 & 5\\
\texttt{lsass-query-by-svchost} & Win & benign & - & Process opened a handle to LSASS & 186/247 & 2\\
\texttt{w32time-runtime-key} & Win & benign & - & Registry value written under a path matching '\textbackslash{}Run' & 300/633 & 3\\
\texttt{gpsvc-service-host} & Win & benign & - & New service host process started as SYSTEM & 174/365 & 2\\
\texttt{linux-shadow-read} & Lin & malicious & T1003.008 & Credential file read by a script run with sudo & 50/76 & 6\\
\texttt{linux-wiper-shred} & Lin & malicious & T1485 & shred overwrote a system directory & 317/672 & 7\\
\texttt{linux-kernel-module} & Lin & malicious & T1547.006 & Kernel module loaded with insmod & 78/142 & 6\\
\texttt{linux-account-created} & Lin & malicious & T1136.001 & Local account created & 76/127 & 5\\
\texttt{linux-sudoers-nopasswd} & Lin & malicious & T1548.003 & Passwordless sudo rule written & 141/245 & 5\\
\texttt{linux-ld-preload-hijack} & Lin & malicious & T1574.006 & Freshly compiled binary executed as root from a home directory & 150/265 & 8\\
\texttt{linux-service-stopped} & Lin & malicious & T1489 & Service state changed by systemctl as root & 1602/3951 & 5\\
\texttt{linux-dpkg-service-start} & Lin & benign & - & Service state changed by systemctl as root & 1602/3951 & 3\\
\texttt{linux-motd-discovery} & Lin & benign & - & Session discovery command run as root & 317/672 & 3\\
\texttt{linux-motd-tmp-cleanup} & Lin & benign & - & Files deleted by rm as root & 317/672 & 3\\
\bottomrule\end{tabular}\end{table*}

